\subsection{Data Infrastructure}
\label{subsec::data_infrastructure}
The Cosmos 3 training corpus is drawn from tens of billions of image and video candidates spanning diverse modalities, domains, and tasks. Operating at this scale demands a data infrastructure that can simultaneously (1) transform raw multimodal data into training-ready samples through large-scale distributed processing, (2) support embedding-based retrieval, clustering, and deduplication, and (3) enable interactive dataset visualization, inspection, and debugging. To meet these requirements, we developed \textbf{SILA} (Scalable Infrastructure for Large-scale data processing and Annotation), a scalable multimodal data infrastructure platform that consolidates storage, metadata management, distributed processing, semantic retrieval, and dataset visualization into a single extensible framework for large-scale data curation and management.

SILA is built around a clean separation between pipeline logic and infrastructure mechanics. Researchers declare typed processing stages---specifying the columns each stage consumes and the outputs it produces---while the platform transparently handles dataset sharding, distributed execution, fault tolerance, checkpointing, metadata updates, and asset registration. This abstraction makes it straightforward to incorporate new data sources, foundation models, and processing stages, including filtering, captioning, embedding generation, scoring, and tagging, without requiring researchers to develop distributed-systems expertise. The result is a platform that lets curation evolve at the pace of research: new signals can be added, recomputed, or replaced incrementally as models, quality criteria, and training recipes change.

\subsubsection{Large-Scale Data Processing}
Multimodal data curation is an iterative enrichment process rather than a single offline preprocessing pass. Raw text, image, and video samples are repeatedly transformed, filtered, annotated, and reprocessed as models, quality criteria, labels, and training recipes evolve. Scaling this workflow is challenging because the pipeline is both low-yield and highly iterative: only a small fraction of raw candidates ultimately survive into training, meaning that inefficient scans, copies, or model inference are disproportionately spent on samples that are later discarded. At the same time, stages such as ingestion, splitting and transcoding, embedding generation, deduplication, filtering, taxonomy tagging, captioning, and sharding repeatedly operate over the same samples. Supporting this workflow therefore requires efficient mechanisms for repeated transformations and incremental recomputation across billions of multimodal samples.

These challenges become even more pronounced under distributed execution. Curation workloads must run continuously on shared clusters with fragmented and dynamically changing GPU availability rather than assuming a single monolithic allocation. The infrastructure must coordinate many distributed workers, avoid duplicate computation, recover from failures, manage heterogeneous CPU- and GPU-bound stages, and continue processing unfinished work as resources become available. To support this execution model, SILA combines a unified data layer with fragment-level coordination and fault recovery, staged distributed execution, node-local model serving, opportunistic cluster utilization, and agent-friendly operational interfaces.

\paragraph{Unified data layer.}
SILA organizes data curation as a unified columnar Lance dataset~\citep{pace2025lance}, where each row represents a data sample and each typed column represents a curation signal such as a caption, tag, quality score, or annotation. This replaces the legacy table-per-pipeline architecture used in earlier infrastructure, Cosmos-Predict 1.0~\citep{cosmos_v1} and 2.5~\citep{cosmos_predict2p5}, where each pipeline wrote to its own Postgres table and outputs were later synchronized into Databricks through Change Data Capture (CDC). As the number of pipelines and metadata fields grew, the table-per-pipeline design required increasingly complex joins across large tables to reconstruct the state of a single sample. These joins became expensive at scale and made even simple operational queries difficult to express without detailed knowledge of join keys, table relationships, and pipeline-specific schemas. In contrast, SILA incrementally enriches the same logical sample by appending new typed columns to a shared Lance table. This unified representation naturally matches multimodal curation workloads, where most stages augment existing samples with additional metadata rather than creating new entities. By co-locating dataset contents, metadata, and processing state, the system can efficiently support large-scale scans, point lookups, incremental recomputation, and discovery of unfinished work directly from Lance fragment metadata without expensive startup joins.

\paragraph{Fragment-level coordination and fault recovery.}
Curation at this scale runs at high concurrency: many distributed workers within a single job, and many independent jobs in parallel, read from and write to the same continuously evolving Lance dataset. Without explicit coordination, workers must rely on expensive startup queries, randomized sampling, or post-hoc filtering of already processed samples to avoid overlap. These approaches delay job startup, permit duplicate work, and complicate recovery when long-running jobs are preempted or interrupted. SILA instead coordinates distributed curation directly at the Lance-fragment level. Workers discover unfinished fragments from Lance metadata and acquire time-limited leases before processing them, while the Lance dataset itself remains the source of truth for completion state. Lease ownership is maintained through periodic heartbeats; when heartbeats stop, the lease expires and another worker can reclaim the fragment, enabling automatic recovery from failures, preemption, or endpoint crashes without manual cleanup. Because a single fragment may contain many samples and require hours of model inference, SILA further partitions claimed fragments into smaller processing segments, writes completed segments as durable checkpoints, and atomically commits the full fragment back into the Lance through a single metadata update once all segments finish. This decouples the recovery unit from the visibility unit: interrupted jobs resume from completed segments while downstream readers observe only fully committed fragment outputs.

\paragraph{Staged Ray execution.}
Curation pipelines combine heterogeneous operations with very different resource profiles, including data loading, decoding, model inference, postprocessing, writing, and committing. If these operations are executed through a single undifferentiated control loop, fast upstream stages can accumulate intermediate outputs while slower inference or commit stages become bottlenecks. SILA instead executes each claimed fragment through a staged Ray pipeline engine~\citep{moritz2018raydistributedframeworkemerging}. Framework-managed stages handle loading, writing, and committing, while user-defined preprocess, compute, and postprocess stages execute in separate Ray actor pools with independently configured worker counts and resource requirements. Backpressure limits in-flight work across stage boundaries, preventing fast I/O-heavy stages from overwhelming slower downstream stages and forcing Ray object-store spill to disk.

\paragraph{Node-local model endpoints.}
Foundation-model curation workloads often require serving large captioning, embedding, tagging, or scoring models while many pipeline workers concurrently process data. Centralized inference services can become bottlenecks at scale, while requiring one large contiguous GPU allocation reduces the ability to exploit fragmented cluster availability. SILA instead launches node-local model-serving endpoints using systems such as vLLM~\citep{vllm} and passes local endpoint information directly to stage workers. Workers then invoke node-local services for inference, allowing model-heavy curation stages to scale across available nodes while decoupling data-parallel pipeline execution from model-serving placement.

\paragraph{Opportunistic cluster utilization.}
Large-scale curation must run continuously alongside training workloads on shared clusters, where GPU availability is often fragmented and dynamically changing. Instead of requiring one large monolithic allocation, SILA decomposes curation into fine-grained distributed jobs that can execute incrementally as resources become available. The system supports execution backends such as DGX Cloud Lepton~\citep{nvidia2026dgxcloudlepton} and Slurm~\citep{jette2023slurmarchitecture}, allowing pipelines to opportunistically utilize idle or partially available GPU capacity. This improves cluster utilization and enables continuous data processing without requiring large contiguous GPU reservations.

\paragraph{Agentic job orchestration.}
As AI agents become increasingly capable at tool use and long-horizon execution, SILA exposes large-scale curation workflows through agent-friendly operational interfaces, including reusable skills, command-line interfaces (CLIs), and structured job metadata. Long-running orchestration agents periodically monitor curation jobs, inspect logs and execution metadata, relaunch failed stages, and coordinate operational recovery automatically. Through this continuous monitoring loop, the agents can track pipeline progress, identify stalled or unhealthy workers, verify dataset coverage, trigger incremental recomputation when new models or filtering criteria are introduced, and notify engineers about failures, recoveries, and execution status as distributed jobs evolve over time.

Together, these design choices substantially improved the efficiency of large-scale curation. By eliminating expensive startup table joins and replacing randomized work selection with fragment-level discovery and coordination, SILA reduced job startup latency from 30--60 minutes to roughly 5 minutes, depending on the stage and model configuration. Combined with staged execution, checkpointing, and improved cluster utilization, the new infrastructure achieved a $10\times$ throughput increase over the previous architecture. In peak production windows, individual SILA stages processed billions of row-level annotations per day, reducing large captioning and curation campaigns from month-scale operations to week-scale iteration cycles.

Beyond improving scalability and throughput, SILA also simplifies pipeline development by hiding much of the operational complexity behind a dataset-centric interface. Pipeline authors specify only the input columns they consume and the output fields they produce, while the framework handles schema registration, column creation, fragment discovery, work coordination, checkpointing, and committing results back to the shared dataset. As a result, adding new captioning, scoring, tagging, or filtering stages no longer requires creating new storage tables, writing synchronization logic, or manually coordinating distributed workers. Researchers can instead iterate by incrementally adding new typed columns, reusing previously computed outputs, and recomputing only samples whose required fields are missing.

\subsubsection{Embedding Storage and Semantic Retrieval}

Semantic retrieval workloads require both vector similarity search and metadata-aware filtering over billions of multimodal data. In the previous architecture, storing high-dimensional embeddings directly in SQL tables significantly inflated table size, increasing I/O overhead for joins, scans, and operational queries. As a result, embeddings were exported into a separate vector database, while metadata used for pre-filtering, post-filtering, and result interpretation remained in relational storage. However, embeddings, metadata, and filtering criteria evolve continuously during curation, requiring frequent synchronization and migration between two systems whenever new embedding models, metadata fields, or search filters are introduced.

SILA instead stores embeddings directly alongside sample metadata in Lance, allowing LanceDB to build vector indexes over the primary dataset rather than requiring a separate vector database~\citep{lancedb2026vectorindexes}. Because embeddings are stored in Lance data files rather than inline relational rows, large embedding payloads do not inflate metadata tables or slow operational queries. By co-locating embeddings, metadata, and vector indexes within the same storage layer, SILA supports semantic retrieval, clustering, and deduplication directly over the curated dataset while keeping search results consistent with the latest curation state.

In production, SILA performs semantic retrieval, clustering, and deduplication over a 4096-dimensional embedding column covering tens of billions of rows using LanceDB IVF\_PQ indexes with cosine similarity. The deployed Approximate Nearest Neighbor (ANN) configuration uses 64K IVF partitions together with PQ-compressed embeddings to support billion-scale retrieval workloads efficiently. Because the vector indexes are built directly over the primary Lance datasets, semantic retrieval operates over the same storage layer that maintains the latest curation metadata and filtering state. This allows metadata-aware filtering, nearest-neighbor retrieval, and downstream dataset analysis to remain synchronized with continuously evolving embeddings, annotations, and curation outputs without requiring synchronization between separate vector and metadata systems.

\subsubsection{Dataset Visualization, Inspection, and Debugging}

At production scale, data curation must be observable as well as scalable: researchers need to understand pipeline coverage, track how curation outputs evolve over time, and diagnose why individual samples pass or fail quality criteria. SILA therefore treats visualization, inspection, and debugging as first-class components of the data infrastructure. Its tools operate directly over Lance tables, connecting aggregate pipeline progress, representative development subsets, sample-level inspection, and downstream analytical views within the same shared curation substrate.

\begin{itemize}

\item \textit{Development and pipeline validation.}
SILA provides utilities for constructing small development Lance tables from production datasets while preserving the schema and representative operational characteristics of the full corpus. Because these development tables closely mirror production data, the same pipelines can run unchanged in both environments, allowing researchers to validate correctness and execution behavior before launching large-scale production jobs.

\item \textit{Dataset inspection and progress analysis.}
To support large-scale curation monitoring, SILA exposes both aggregate progress analyzers and interactive inspection tools directly over Lance tables. Fragment-level metadata is used to estimate per-column coverage and pipeline completion without requiring full table scans, while interactive viewers allow researchers to sample rows, render media, and inspect the associated captions, scores, annotations, and schema metadata. Because Lance supports efficient random row-level access, these inspection workflows can operate directly over the primary curation tables, avoiding the expensive scans and limited row-level retrieval patterns common in traditional Parquet-based data lakes. This allows researchers to move quickly from pipeline-level progress monitoring to sample-level debugging within the same dataset.

\item \textit{Analytical querying integration.}
For large analytical workloads, SILA supports Online Analytical Processing (OLAP) queries used for large-scale aggregation, reporting, and dashboarding over the curated corpus. These queries are simpler to express because SILA stores each sample and its curation signals in a wide Lance table: captions, tags, scores, embeddings, filtering decisions, and processing state can be selected and filtered from one logical dataset rather than reconstructed through joins across pipeline-specific tables. After selecting the relevant columns and cohorts, analytical backends can compute the required aggregations for coverage reports, quality dashboards, dataset audits, and training-set analysis. SILA keeps the Lance table as the source of truth for curation state, media, metadata, embeddings, vector indexes, and row-level inspection, while treating downstream analytical execution as a deployment choice. In practice, SILA can scan the authoritative Lance tables to materialize query-ready snapshots or projections, including Parquet files~\citep{parquet2013}, and execute OLAP workloads using the available compute backend, such as Databricks, Spark clusters~\citep{zaharia2016apache}, or Slurm-backed batch jobs.

\end{itemize}

Across processing, retrieval, and inspection, SILA closes the loop on the three objectives that define the Cosmos 3 data infrastructure: transforming raw multimodal corpora into training-ready samples, organizing them for semantic retrieval and deduplication, and keeping them inspectable throughout the curation lifecycle. By unifying assets, curation signals, embeddings, vector indexes, and execution state within a single Lance-backed substrate, SILA turns data curation from a sequence of one-off preprocessing jobs into a continuously evolving production workflow. New models and quality criteria can be applied incrementally, distributed enrichment stages can recover and scale across shared heterogeneous clusters, and researchers can move seamlessly from corpus-level progress monitoring to sample-level debugging. This integrated workflow enables the Cosmos 3 training corpus to scale to tens of billions of multimodal candidates while remaining searchable, auditable, and continuously improvable.
